\begin{textsample}{POS Dim 1 – pe – Score 29.00 – pe/pe\_ef\_26.txt}  \label{ex:f1_pos_189}
3.3.7 A LÍNGUA PORTUGUESA NO ENSINO FUNDAMENTAL : \textbf{UM} VIÉS \textbf{METODOLÓGICO} Conforme já sinalizado neste documento, o ensino da Língua Portuguesa parte do princípio de \textbf{que} a língua é \textbf{um} objeto de interação social.
Essa concepção se concretiza na sala de aula através de sequências didáticas \textbf{que} favoreçam atividades de leitura, voltadas para a interação de autores/leitores, buscando a construção de sentidos de textos lidos/ouvidos e/ou sinalizados, atividades de produção de textos orais e escritos, compreendidos \textbf{enquanto} propostas de produção de sentidos previamente definidas, em \textbf{que} alguém \textbf{diz} alguma coisa para alguém.
Tomando por base esse viés sócio interacionista, as atividades planejadas, especialmente nos anos iniciais do ensino fundamental, para a apropriação do sistema de escrita, para a \textbf{sistematização} de conhecimentos linguísticos e gramaticais partem de práticas sociais efetivas de leitura, da escrita e da oralidade, para \textbf{que} os estudantes compreendam \textbf{que} a língua \textbf{que} estão estudando é a mesma \textbf{que} usada nas diversas esferas sociais onde eles circulam e interagem.
Nesse contexto, vale \textbf{salientar}, \textbf{que} os estudantes \textbf{que} ingressam no ensino fundamental, são crianças de 06 anos, oriundas da educação infantil, \textbf{que} tem como eixos estruturadores as interações e as brincadeiras.
Conforme orienta o Currículo de Pernambuco, à luz da BNCC, é \textbf{imprescindível} promover a integração entre cada etapa da educação básica, para \textbf{que} não existam rupturas entre as mesmas.
Essas crianças, embora estejam inseridas socialmente em contextos de letramento e tenham vivenciado as múltiplas linguagens, encontram muitos desafios na apropriação de conhecimentos relativos à língua.
Nesse sentido, atividades planejadas para a apropriação do sistema de escrita e a consolidação do processo de alfabetização e letramento podem ter na dimensão lúdica \textbf{ancorada}, sobretudo, nos eixos estruturantes e integradores \textbf{uma} possibilidade de atenuar e superar tais desafios.
Atividades de jogos e brincadeiras, além de minimizar essa ruptura, fortalecem a ação pedagógica em dois aspectos : a ) promovem aprendizagens diferenciadas, porque possibilitam ao professor dinamizar suas aulas, estimulando o interesse e a participação dos estudantes, tornando as aulas mais atraentes e \textbf{interessantes} ; b ) criam possibilidades de aprendizagem dinâmica e de qualidade, \textbf{uma} vez \textbf{que} a criança terá a oportunidade de aprender brincando, possibilitando o ensino prazeroso e a aprendizagem significativa, \textbf{condizente} com a sua realidade.
Isso é possível porque estão, nas brincadeiras, as formas mais autênticas de se relacionar, de interagir com o outro e de se apropriar do mundo, dos conhecimentos.
Portanto, o \textbf{estímulo} encontrado no lúdico estabelece a integração dos conteúdos escolares ao cotidiano dos estudantes.
Considerando \textbf{que} as crianças, em processo de alfabetização, precisam entender e usar a representação fonética e gráfica de muitos caracteres/sons e usá -los em palavras e textos do seu cotidiano, a perspectiva do lúdico torna -se \textbf{um} facilitador desse processo ensino-aprendizagem, dada a complexidade de abstração e representação de tais conhecimentos.
Brincadeiras/jogos, como por exemplo, bingo de letras, bingos de sons, cruzadinhas, jogos de rimas, dominós diversos, jogos de análise fonológica, são importantes instrumentos para o aprendizado da língua, porque levam os estudantes a pensar nas palavras em sua dimensão não só sonora, gráfica, mas também semântica.
Portanto, jogos e brincadeiras devem ser usados para fazer com \textbf{que} a criança reflita, ordene, desorganize, desmonte e remonte o mundo a sua maneira.
Na ação de brincar, ela tem oportunidade de aprender regras, normas, valores e também conteúdos \textbf{conceituais}, atitudinais e procedimentais nas mais diversas formas de conhecimento.
O lúdico também favorece a autoestima e a interação entre os pares, propiciando situações de aprendizagem e desenvolvimento das capacidades cognitivas e afetivas das crianças em sala de aula.
Através da brincadeira e de jogos tanto impressos quanto digitais, as crianças vão construindo seu aprendizado, entretanto, essas atividades devem ser planejadas, diversificadas, e ter como objetivo principal o amplo desenvolvimento da criança, em todas as dimensões humanas, sejam essas biológicas, motoras, cognitivas, afetivas, social, ética, entre outras.
Por outro \textbf{lado}, embora a inserção de atividades lúdicas nos anos finais do ensino fundamental esteja presente em muitas estratégias didáticas e seja, sobretudo, significativa para os processos de interação e socialização, contribuindo para a ampliação e o aprofundamento das habilidades e competências já consolidadas pelos estudantes ou em via de consolidação como também para \textbf{aproximar} e fortalecer os vínculos ( entre ) professor e estudantes, faz -se necessário, nesta etapa da educação básica, desenvolver metodologias, práticas e ações pedagógicas \textbf{que} visem à promoção e ao fortalecimento da autonomia, do caráter crítico e investigativo, da fluência e, especialmente, do protagonismo sem \textbf{perder} o foco na formação de leitores e de produtores de textos ( orais e escritos ) aptos a exercerem sua cidadania e \textbf{humanidade} plena e conscientemente.
Pela natureza dos objetos de conhecimentos e habilidades selecionados para os anos finais, e ainda pelo contexto social em \textbf{que} os estudantes estão inseridos ( de alta modernidade ), professores e estudantes precisam conceber os processos de ensino numa perspectiva colaborativa, interdisciplinar, desvinculada da \textbf{memorização} e transmissão de conteúdos, portanto, centrada nas aprendizagens e na busca de novas formas de interação entre todos os \textbf{atores} sociais ( estudantes, professores, familiares, comunidade etc. ).
Nessa \textbf{direção}, as \textbf{pedagogias} da autonomia[6 ] e do protagonismo[7 ] se empenharão em garantir diálogos entre as culturas locais dos estudantes ( presentes na escola ) e a cultura valorizada por ela, entre as diferentes linguagens e modos de representar a realidade, entre os diferentes papeis \textbf{exigidos} por \textbf{uma} sociedade altamente tecnológica, midiática e semiótica, contribuindo para a formação de \textbf{sujeitos}, primeiramente, \textbf{corresponsáveis} por suas aprendizagens e desenvolvimento e, consequentemente, autônomos, solidários, participativos e éticos.
A Língua Portuguesa, nesse contexto, através das práticas de linguagem juntamente com diferentes estratégias, procedimentos e condições de produção, efeitos de sentido, de articulação entre os recursos linguísticos e semióticos etc., possibilita a construção e expressão das diferentes identidades, para inserção e participação social, para o projeto de vida, para o fortalecimento de posturas e espaços democráticos, para leitura crítica e plurissignificativa da realidade e dos bens simbólicos e culturais, entre tantos outros aspectos, \textbf{uma} vez \textbf{que}, segundo Antunes ( 2004 ), a língua se atualiza a serviço da comunicação intersubjetiva, em situações de atuação social e através de práticas discursivas, materializadas em textos. [ 6 ] : Prática \textbf{que} respeita, na visão freireana, a cultura e os conhecimentos empíricos e subjetivos dos estudantes, \textbf{visto} \textbf{que} são sujeitos sócio-histórico-culturais, visando a \textbf{uma} melhor integração entre o SER e o SABER dos educandos.
Atrelada à ética e à \textbf{dignidade} do ser, a autonomia do estudante respalda -se na concepção de \textbf{que} “ formar é muito mais do \textbf{que} treinar o \textbf{educando} no desempenho de destreza ”. ( \textbf{FREIRE}, 1997, p.3 ). [ 7 ] : Concebida como ação educativa \textbf{que} possibilita aos estudantes se envolverem em atividades direcionadas a soluções de problemas reais com liberdade e \textbf{compromisso}, tornando -se personagens principais nas soluções de problemas da escola, da comunidade e até da sociedade.
Por sua vez, essa atuação coopera para a construção da autonomia, autoconfiança e autodeterminação, culminando na construção de identidades e do projeto de vida ( COSTA, 2018 ).

% matched lemmas: ancorar, aproximar, ator, compromisso, conceitual, condizente, corresponsável, dignidade, direção, dizer, educar, enquanto, estímulo, exigir, freire, humanidade, imprescindível, interessante, lado, memorização, metodológico, pedagogia, perder, que, salientar, sistematização, sujeito, um, ver
\end{textsample}
